\chapter{Cho em cơ hội}
\markboth{Cho em cơ hội}{Give Another Chance}


\section{Làm lại không đáng xấu hổ — bỏ luôn mới đáng tiếc.}
\begin{truyen}{Làm lại không đáng xấu hổ — bỏ luôn mới đáng tiếc.}{A second chance matters.}
\chuhoa{C}{ó học sinh nản chỉ vì nghĩ mình đã sai một lần là hết đường. Nhưng một cánh cửa làm lại mở ra đúng lúc sẽ cứu được rất nhiều điều còn dang dở.}
\textit{(Some students are discouraged simply because they think one mistake ends everything. But a door to try again, opened at the right time, can save many unfinished things.)}
\end{truyen}

\begin{baihoc}
Cơ hội làm lại không dung túng sự yếu kém; nó trao quyền cho sự trưởng thành.
\textit{(A second chance does not excuse weakness; it gives growth another chance.)}
\end{baihoc}

\begin{ghinhoanh}
\textbf{Short English to remember:}

Try again. That is not shame.

\end{ghinhoanh}

\begin{tuhocnhanh}
1. Em đang muốn làm lại điều gì? \textit{(What do you wish to try again?)}

2. Điều gì làm em ngại xin thêm một cơ hội? \textit{(What makes you hesitate to ask for another chance?)}
\end{tuhocnhanh}
\ngancach


\section{Sai một bài không có nghĩa em sai cả con đường.}
\begin{truyen}{Sai một bài không có nghĩa em sai cả con đường.}{One failure is not all.}
\chuhoa{N}{hiều em kéo một điểm kém thành kết luận về bản thân. Nhưng một bài sai chỉ là một đoạn nhỏ, không phải bản án cho cả chặng học tập.}
\textit{(Many students turn one low mark into a conclusion about themselves. But one wrong test is only a small part, not a verdict on the whole journey.)}
\end{truyen}

\begin{baihoc}
Đừng biến một thất bại nhỏ thành danh tính của mình.
\textit{(Do not turn one small failure into your identity.)}
\end{baihoc}

\begin{ghinhoanh}
\textbf{Short English to remember:}

One bad test is not your whole story.

\end{ghinhoanh}

\begin{tuhocnhanh}
1. Em có đang đánh giá mình từ một lần sai không? \textit{(Are you judging yourself from one failure?)}

2. Nếu nhìn rộng hơn, em còn điều gì làm được? \textit{(If you look wider, what else have you done well?)}
\end{tuhocnhanh}
\ngancach


\section{Hôm nay chưa được thì mai làm lại — miễn là em còn ở lại.}
\begin{truyen}{Hôm nay chưa được thì mai làm lại — miễn là em còn ở lại.}{Stay in the game.}
\chuhoa{Đ}{iều quan trọng nhất khi học sinh nản không phải ép các em thắng ngay, mà giữ các em ở lại với hành trình đủ lâu để có ngày làm được.}
\textit{(The most important thing when students are discouraged is not forcing them to win immediately, but keeping them in the journey long enough for success to arrive.)}
\end{truyen}

\begin{baihoc}
Còn ở lại là còn cơ hội đổi kết quả.
\textit{(As long as you stay, you can still change the outcome.)}
\end{baihoc}

\begin{ghinhoanh}
\textbf{Short English to remember:}

Stay. Tomorrow can change.

\end{ghinhoanh}

\begin{tuhocnhanh}
1. Điều gì làm em muốn bỏ giữa chừng nhất? \textit{(What most makes you want to quit halfway?)}

2. Điều gì giúp em ở lại thêm một chút? \textit{(What helps you stay a little longer?)}
\end{tuhocnhanh}
\ngancach


\section{Thầy sẵn sàng cho em thêm một lần — nếu em cũng sẵn sàng với chính mình.}
\begin{truyen}{Thầy sẵn sàng cho em thêm một lần — nếu em cũng sẵn sàng...}{Use the chance well.}
\chuhoa{C}{ơ hội chỉ thật sự có nghĩa khi học sinh cũng chịu bước vào lại với thái độ mới. Thầy mở cửa, nhưng chính em phải bước qua ngưỡng đó.}
\textit{(A second chance only matters when the student is also willing to return with a new attitude. The teacher opens the door, but you must step through it.)}
\end{truyen}

\begin{baihoc}
Cơ hội là món quà; trách nhiệm là phần em phải tự mang.
\textit{(Opportunity is a gift; responsibility is the part you carry yourself.)}
\end{baihoc}

\begin{ghinhoanh}
\textbf{Short English to remember:}

Second chances need second effort.

\end{ghinhoanh}

\begin{tuhocnhanh}
1. Nếu được thêm một cơ hội, em sẽ làm khác điều gì? \textit{(If you had another chance, what would you do differently?)}

2. Em đã thật sự sẵn sàng chưa? \textit{(Are you truly ready?)}
\end{tuhocnhanh}
\ngancach


\section{Đi chậm lại cũng được — miễn là em chưa quay lưng.}
\begin{truyen}{Đi chậm lại cũng được — miễn là em chưa quay lưng.}{Slow is still forward.}
\chuhoa{N}{hiều em ngại vì thấy mình chậm hơn bạn. Nhưng đi chậm không đáng sợ bằng quay lưng với việc học chỉ vì xấu hổ.}
\textit{(Many students feel ashamed for moving slower than others. But moving slowly is far less dangerous than turning away from learning out of shame.)}
\end{truyen}

\begin{baihoc}
Chậm không phải thua; bỏ luôn mới là mất.
\textit{(Slow is not losing; walking away is.)}
\end{baihoc}

\begin{ghinhoanh}
\textbf{Short English to remember:}

Slow is okay. Leaving is not.

\end{ghinhoanh}

\begin{tuhocnhanh}
1. Em có đang xấu hổ vì nhịp của mình không? \textit{(Are you ashamed of your pace?)}

2. Em cần tự nói gì để bớt so sánh? \textit{(What do you need to tell yourself to compare less?)}
\end{tuhocnhanh}
\ngancach


\section{Một lỗi được sửa tốt còn quý hơn một bài làm đẹp mà không hiểu.}
\begin{truyen}{Một lỗi được sửa tốt còn quý hơn một bài làm đẹp mà không hiểu.}{Correcting builds mastery.}
\chuhoa{C}{ó học sinh ngại sửa sai vì nghĩ sai là xấu. Nhưng biết sửa cho tử tế mới là dấu hiệu của người học đang lớn lên thật.}
\textit{(Some students hate correcting mistakes because they think mistakes are shameful. But learning to fix them well is a real sign of growth.)}
\end{truyen}

\begin{baihoc}
Sửa sai không làm em nhỏ đi; nó làm em chắc hơn.
\textit{(Correcting mistakes does not make you smaller; it makes you steadier.)}
\end{baihoc}

\begin{ghinhoanh}
\textbf{Short English to remember:}

Fixing mistakes is real learning.

\end{ghinhoanh}

\begin{tuhocnhanh}
1. Em thường né hay nhìn lại lỗi của mình? \textit{(Do you avoid or revisit your mistakes?)}

2. Lần sửa sai nào khiến em tiến bộ rõ nhất? \textit{(Which correction helped you improve the most?)}
\end{tuhocnhanh}
\ngancach


\section{Thầy không giữ lỗi cũ để nhắc mãi — thầy giữ hy vọng để chờ em đổi.}
\begin{truyen}{Thầy không giữ lỗi cũ để nhắc mãi — thầy giữ hy vọng để c...}{Hope over labels.}
\chuhoa{N}{hãn dán cũ làm học sinh tưởng mình không thể khác đi. Hy vọng của người thầy mới là điều giữ cho các em còn muốn thử lại.}
\textit{(Old labels make students believe they cannot change. A teacher's hope is what keeps them willing to try again.)}
\end{truyen}

\begin{baihoc}
Người thầy không nhốt học sinh trong lỗi cũ.
\textit{(A true teacher does not trap students inside old mistakes.)}
\end{baihoc}

\begin{ghinhoanh}
\textbf{Short English to remember:}

Hope changes more than labels.

\end{ghinhoanh}

\begin{tuhocnhanh}
1. Có nhãn nào em bị gắn mà em muốn gỡ bỏ? \textit{(Is there any label you want to remove from yourself?)}

2. Em muốn người khác hy vọng điều gì ở mình? \textit{(What do you want others to keep hoping about you?)}
\end{tuhocnhanh}
\ngancach


\section{Em chưa làm được hôm nay không có nghĩa em không làm được ngày mai.}
\begin{truyen}{Em chưa làm được hôm nay không có nghĩa em không làm được ...}{Not yet, not never.}
\chuhoa{C}{âu nói khác nhau chỉ ở một chữ `chưa', nhưng chính chữ đó cứu học sinh khỏi cảm giác bế tắc. `Chưa' còn mở đường; `không thể' thì đóng sập lại.}
\textit{(The difference may be only the word `yet', but that word saves students from hopelessness. `Not yet' keeps the road open; `impossible' shuts it down.)}
\end{truyen}

\begin{baihoc}
Đừng biến `chưa' thành `không bao giờ'.
\textit{(Do not turn `not yet' into `never'.)}
\end{baihoc}

\begin{ghinhoanh}
\textbf{Short English to remember:}

Not yet is still hope.

\end{ghinhoanh}

\begin{tuhocnhanh}
1. Việc gì em đang nghĩ mình `không thể'? \textit{(What are you thinking you `cannot do'?)}

2. Nếu đổi thành `chưa', em thấy khác thế nào? \textit{(How does it feel if you change that to `not yet'?)}
\end{tuhocnhanh}
\ngancach


\section{Thử thêm một lần không làm em yếu — nó làm em bền.}
\begin{truyen}{Thử thêm một lần không làm em yếu — nó làm em bền.}{Persistence builds strength.}
\chuhoa{N}{gười ta thường khen kết quả, nhưng sự bền mới là thứ đem kết quả tới. Một lần thử nữa đôi khi là khác biệt giữa dừng lại và trưởng thành.}
\textit{(People usually praise results, but endurance is what brings results. One more try is sometimes the difference between stopping and growing up.)}
\end{truyen}

\begin{baihoc}
Kiên trì là sức mạnh thường bị đánh giá thấp.
\textit{(Persistence is a strength people often underestimate.)}
\end{baihoc}

\begin{ghinhoanh}
\textbf{Short English to remember:}

One more try builds grit.

\end{ghinhoanh}

\begin{tuhocnhanh}
1. Em đã từ bỏ điều gì quá sớm? \textit{(What have you given up too early?)}

2. Một lần thử nữa hôm nay có thể là gì? \textit{(What could one more try be today?)}
\end{tuhocnhanh}
\ngancach


\section{Nếu em quay lại, thầy vẫn mở cửa.}
\begin{truyen}{Nếu em quay lại, thầy vẫn mở cửa.}{The door is still open.}
\chuhoa{C}{ó học sinh rời khỏi việc học bằng im lặng. Điều các em cần nghe không phải là trách cứ, mà là biết cánh cửa quay lại vẫn còn đó.}
\textit{(Some students drift away from learning in silence. What they need to hear is not blame, but that the door back is still there.)}
\end{truyen}

\begin{baihoc}
Một cánh cửa mở có thể kéo học sinh trở lại đúng lúc nhất.
\textit{(An open door can bring students back at the exact right moment.)}
\end{baihoc}

\begin{ghinhoanh}
\textbf{Short English to remember:}

The door back matters.

\end{ghinhoanh}

\begin{tuhocnhanh}
1. Nếu lỡ rời khỏi một việc tốt, em có dám quay lại không? \textit{(If you drift away from something good, do you dare to return?)}

2. Điều gì sẽ giúp em quay lại dễ hơn? \textit{(What would make returning easier for you?)}
\end{tuhocnhanh}
\ngancach